% (Ch20 misconception items are addressed in the status paragraphs of prop:heavy-qcdf-lp,
%  thm:heavy-nrqcd-decay and conj:heavy-nrqcd-production.)
\paragraph{本章小结。}
$B\to D\pi$ 与大反冲下的重到轻形状因子有 SCET 中的全阶论证，前提是模式完备性与端点压低；无粲两体衰变的 BBNS 公式在 $\order{\as}$ 已证明、$T^{\rm I}$ 到 NNLO 显式红外有限，全阶证明尚无共识（[P]），其幂次修正因卷积发散不能按标准共线形式因子化（[V]），$X_{A,H}$ 只是模型。单举 $B$ 衰变率没有 $1/m_b$ 项，谱的端点区是硬函数、喷注函数与形状函数的卷积。NRQCD 衰变公式在微扰论中有全阶论证，匹配系数 $c_v$ 到单圈严格（[R]）；产生的 NRQCD 因子化是猜想，在 NNLO 需要规范完备的 Wilson 线。用于固定运动学点上的强子衰变率，以上结果都还需要局域对偶（[C]）。
